\documentclass[11pt,a4paper]{article}
\usepackage[utf8]{inputenc}
\usepackage[T1]{fontenc}
\usepackage[brazilian]{babel}
\usepackage{lmodern}
\usepackage{amsmath, amsfonts, amssymb}
\usepackage{geometry}
\usepackage{graphicx}
\usepackage{subcaption}
\usepackage{float}
\usepackage{enumitem}
\usepackage{booktabs}
\usepackage{array}
\usepackage{indentfirst}
\usepackage{microtype}
\usepackage{setspace}
\usepackage[font=small,labelfont=bf,labelsep=period]{caption}
\usepackage{hyperref}

\geometry{left=2.2cm, right=2.2cm, top=2.2cm, bottom=2.2cm}
\sloppy
\setstretch{1.14}
\setlength{\textfloatsep}{10pt plus 2pt minus 2pt}
\setlength{\floatsep}{8pt plus 2pt minus 2pt}
\setlength{\abovedisplayskip}{6pt plus 2pt minus 2pt}
\setlength{\belowdisplayskip}{6pt plus 2pt minus 2pt}

\hypersetup{
    colorlinks=false,
    pdfborder={0 0 0},
    pdftitle={Trabalho 3 - Modelos Probabilisticos - Grupo 2 (Star Trek)},
    pdfauthor={Bruno Mariano, Jaco Medeiros Lima, Joao Victor Pariz, Lucas Quiuqui, Marcos Vinicius, Renato Felix}
}

\setlength{\parskip}{3pt}
\setlength{\parindent}{1.25cm}

\begin{document}

% ========================================================
% CABEÇALHO DO RELATÓRIO
% ========================================================
\begin{center}
    {\large \textbf{INSTITUTO MILITAR DE ENGENHARIA}}\\[1mm]
    {\normalsize \textbf{SEÇÃO DE ENGENHARIA ELÉTRICA E DE SISTEMAS (SE/3)}}\\[1mm]
    {\normalsize \textbf{DISCIPLINA: MODELOS PROBABILÍSTICOS (EE-03247)}}\\[3mm]
    \rule{\textwidth}{0.8pt}\\[2.5mm]
    {\Large \textbf{TRABALHO 3}}\\[1.5mm]
    {\large \textbf{ANÁLISE ESTOCÁSTICA DE SINAIS DE VÍDEO CODIFICADOS}}\\[1.5mm]
    {\normalsize \textbf{Grupo 2 — Sequência \textit{Star Trek}}}\\[2.5mm]
    {\normalsize \textbf{Integrantes do Grupo:}}\\[1mm]
    {\normalsize Bruno Mariano \quad Jacó Medeiros Lima \quad João Victor Pariz}\\[1mm]
    {\normalsize Lucas Quiuqui \quad Marcos Vinícius \quad Renato Felix}\\[2mm]
    {\normalsize \textbf{Professora:} Rosângela Fernandes Coelho}\\[1mm]
    {\small Rio de Janeiro, 12 de novembro de 2026}\\[1.5mm]
    \rule{\textwidth}{0.8pt}
\end{center}

\vspace{1mm}

% ========================================================
% INTRODUÇÃO
% ========================================================
\section*{Introdução}

Este relatório apresenta o Trabalho 3 de Modelos Probabilísticos: a análise estocástica de dois sinais de vídeo codificados do filme \textit{Star Trek} (Grupo 2). Os sinais são o número de bits por quadro de cada codificação:
\begin{itemize}[leftmargin=*, nosep]
    \item \texttt{startrek\_h263\_64k.txt} ($X[n]$): codificado em H.263 a $64\text{ kbps}$, com $N_X = 26.192$ amostras;
    \item \texttt{startrek\_mpeg4\_hq.txt} ($Y[n]$): codificado em MPEG-4 de alta qualidade, com $N_Y = 89.998$ amostras e GoP (\textit{Group of Pictures}) de $12$ quadros (\texttt{IBBPBBPBBPBB}).
\end{itemize}

A teoria usada vem das aulas da professora Rosângela e dos Capítulos 6 e 7 do livro de Albuquerque, Fortes e Finamore (J.J.W.) \cite{JJW}.

% ========================================================
% ITEM 1
% ========================================================
\section{Sinais Temporais, Histogramas e Caracterização Gaussiana}

\noindent\textbf{Enunciado (Item 1):} \textit{Plotar os sinais e seus Histogramas. Verificar se podem ser caracterizados (PDF) por um vetor Gaussiano.}

\medskip
Pela Definição 7.23 do J.J.W. \cite{JJW}, um processo discreto é caracterizado por um vetor Gaussiano se todo vetor $\underline{X} = [X[n_1], \dots, X[n_k]]^T$ for conjuntamente Gaussiano. Para isso, cada marginal precisa ser Gaussiana:
\begin{equation}
P_X(x) = \frac{1}{\sqrt{2\pi\sigma_X^2}} \exp\!\left(-\frac{(x - m_X)^2}{2\sigma_X^2}\right), \qquad -\infty < x < \infty
\end{equation}

A Figura \ref{fig:sinais_tempo} mostra os sinais no tempo e a Figura \ref{fig:histogramas} mostra os histogramas junto da Gaussiana de mesma média e desvio-padrão.

\begin{figure}[!htbp]
    \centering
    \begin{subfigure}[b]{0.48\textwidth}
        \centering
        \includegraphics[width=\linewidth]{Programas/resultados/passo1_sinais_e_histogramas/sinal_startrek_h263_64k.png}
        \caption{Sinal $X[n]$ (\texttt{h263\_64k}, $N_X=26.192$).}
    \end{subfigure}
    \hfill
    \begin{subfigure}[b]{0.48\textwidth}
        \centering
        \includegraphics[width=\linewidth]{Programas/resultados/passo1_sinais_e_histogramas/sinal_startrek_mpeg4_hq.png}
        \caption{Sinal $Y[n]$ (\texttt{mpeg4\_hq}, $N_Y=89.998$).}
    \end{subfigure}
    \caption{Número de bits por quadro de cada sinal.}
    \label{fig:sinais_tempo}
\end{figure}

\begin{figure}[!htbp]
    \centering
    \begin{subfigure}[b]{0.48\textwidth}
        \centering
        \includegraphics[width=\linewidth]{Programas/resultados/passo1_sinais_e_histogramas/histograma_gaussiana_startrek_h263_64k.png}
        \caption{Histograma e Gaussiana teórica ($X[n]$).}
    \end{subfigure}
    \hfill
    \begin{subfigure}[b]{0.48\textwidth}
        \centering
        \includegraphics[width=\linewidth]{Programas/resultados/passo1_sinais_e_histogramas/histograma_gaussiana_startrek_mpeg4_hq.png}
        \caption{Histograma e Gaussiana teórica ($Y[n]$).}
    \end{subfigure}
    \caption{Histogramas e f.d.p.\ Gaussiana de mesma média e variância.}
    \label{fig:histogramas}
\end{figure}

Os sinais \textbf{não podem ser caracterizados por um vetor Gaussiano}, pelos motivos abaixo:
\begin{enumerate}[leftmargin=*, nosep]
    \item \textbf{Valores não negativos ($X[n] \ge 0$):} o número de bits nunca é negativo, mas a Gaussiana dá probabilidade a valores negativos. Pela função $Q(\alpha)$, $P(Y < 0) = Q(m_Y/\sigma_Y) = Q(1{,}3638) \approx 8{,}63\%$ no MPEG-4 e $P(X < 0) = Q(m_X/\sigma_X) = Q(2{,}1806) \approx 1{,}46\%$ no H.263.
    \item \textbf{H.263 com dois picos:} o histograma (Figura \ref{fig:histogramas}a) tem dois máximos, em torno de $750$ e $1.450\text{ bits}$, enquanto a Gaussiana tem um só.
    \item \textbf{MPEG-4 assimétrico:} o histograma (Figura \ref{fig:histogramas}b) concentra-se em valores baixos (quadros P e B) e tem cauda longa à direita, até $12.500\text{ bits}$ (quadros I).
\end{enumerate}
Como as marginais não são Gaussianas, o vetor não é Gaussiano.

% ========================================================
% ITEM 2
% ========================================================
\section{Valores Esperados: Média, Variância, Desvio-Padrão e VMQ}

\noindent\textbf{Enunciado (Item 2):} \textit{Valores esperados: média, variância ou desvio-padrão e VMQ entre as duas sequências amostrais (sinais de vídeo codificados).}

\medskip
Os momentos foram estimados pelas médias amostrais:
\begin{align}
m_X &= \frac{1}{N}\sum_{n=0}^{N-1} x[n], & \text{VMQ}_X &= \frac{1}{N}\sum_{n=0}^{N-1} x^2[n], \\
\sigma_X^2 &= \text{VMQ}_X - m_X^2, & \sigma_X &= \sqrt{\sigma_X^2}
\end{align}

O H.263 tem $26.192$ amostras e o MPEG-4 tem $89.998$. Para calcular os momentos conjuntos e a correlação cruzada, os dois sinais precisam ter o mesmo tamanho. Por isso cortou-se do MPEG-4 um bloco contínuo de $26.192$ amostras (índices $k = 8.992$ a $35.184$), escolhido para manter a média e o desvio-padrão da série completa (erro menor que $1\%$). A Figura \ref{fig:corte_mpeg4} compara a f.d.p.\ do sinal original com a do bloco. Os resultados estão nas Tabelas \ref{tab:momentos} e \ref{tab:conjuntos}.

\begin{figure}[!htbp]
    \centering
    \includegraphics[width=0.58\linewidth]{Programas/resultados/passo3_facs/comparacao_pdf_mpeg4_cortada.png}
    \caption{F.d.p.\ do MPEG-4 original ($89.998$ amostras) e do bloco cortado ($26.192$ amostras, $k = 8.992$ a $35.184$).}
    \label{fig:corte_mpeg4}
\end{figure}

\begin{table}[!htbp]
\centering
\caption{Valores esperados dos sinais analisados (\textit{Star Trek}).}
\label{tab:momentos}
\renewcommand{\arraystretch}{1.2}
\small
\begin{tabular}{lrrr}
\toprule
Grandeza & \multicolumn{1}{c}{$X[n]$ (H.263)} & \multicolumn{1}{c}{$Y[n]$ (MPEG-4 completo)} & \multicolumn{1}{c}{$Y_{\text{bloco}}[n]$ (MPEG-4 bloco)} \\
\midrule
Número de amostras ($N$) & $26.192$ & $89.998$ & $26.192$ \\
Média $m$ (bits) & $1.099{,}28$ & $1.644{,}81$ & $1.640{,}61$ \\
Variância $\sigma^2$ (bits$^2$) & $254.130{,}08$ & $1.454.519{,}50$ & $1.481.147{,}70$ \\
Desvio-padrão $\sigma$ (bits) & $504{,}11$ & $1.206{,}03$ & $1.217{,}02$ \\
VMQ $=E[(\cdot)^2]$ (bits$^2$) & $1.462.546{,}37$ & $4.159.911{,}56$ & $4.172.745{,}19$ \\
Valor eficaz $\sqrt{\text{VMQ}}$ (bits) & $1.209{,}36$ & $2.039{,}59$ & $2.042{,}73$ \\
\bottomrule
\end{tabular}
\end{table}

\begin{table}[!htbp]
\centering
\caption{Momentos conjuntos entre $X[n]$ e $Y_{\text{bloco}}[n]$ ($N = 26.192$).}
\label{tab:conjuntos}
\renewcommand{\arraystretch}{1.2}
\small
\begin{tabular}{lr}
\toprule
Grandeza & Valor \\
\midrule
Correlação cruzada $R_{XY}[0] = E[XY]$ & $1.800.524{,}75\ \text{bits}^2$ \\
Covariância cruzada $K_{XY}[0] = R_{XY}[0] - m_X m_Y$ & $-2.963{,}91\ \text{bits}^2$ \\
Coeficiente de correlação $\rho_{XY}[0] = K_{XY}[0]/(\sigma_X \sigma_Y)$ & $-0{,}0048$ \\
VMQ da diferença $E[(X - Y)^2]$ & $2.034.243{,}53\ \text{bits}^2$ \\
\bottomrule
\end{tabular}
\end{table}

Confere a relação $\text{VMQ} = \sigma^2 + m^2$:
\begin{align*}
\text{H.263: } \quad & 254.130{,}08 + (1.099{,}28)^2 = 1.462.546{,}37\text{ bits}^2 = \text{VMQ}_X \\
\text{MPEG-4: } \quad & 1.454.519{,}50 + (1.644{,}81)^2 = 4.159.911{,}56\text{ bits}^2 = \text{VMQ}_Y
\end{align*}
O MPEG-4 tem média $50\%$ maior e variância cerca de $5{,}7$ vezes a do H.263. O H.263 trabalha com taxa limitada a $64\text{ kbps}$, então o tamanho dos quadros varia menos.

% ========================================================
% ITEM 3
% ========================================================
\section{Funções de Autocorrelação (FACs)}

\noindent\textbf{Enunciado (Item 3):} \textit{Plotar as 3 FACs.}

\medskip
A FAC é $R_X[k] = E[X[n]X[n+k]]$ e a correlação cruzada é $R_{XY}[k] = E[X[n]Y[n+k]]$. Como a média não é nula:
\begin{itemize}[leftmargin=*, nosep]
    \item Na origem ($k = 0$): $R_X[0] = E[X^2] = \text{VMQ}_X = \sigma_X^2 + m_X^2$;
    \item Para $k$ grande: $R_X[k] \to m_X^2$ (nível DC).
\end{itemize}
Normalizando por $R_X[0] = \text{VMQ}_X$:
\begin{equation}
r_X[k] = \frac{R_X[k]}{R_X[0]} \implies r_X[0] = 1{,}0 \quad \text{e} \quad \lim_{k\to\infty} r_X[k] = \frac{m_X^2}{\text{VMQ}_X}
\end{equation}

As FACs estão na Figura \ref{fig:facs_individuais} e a correlação cruzada na Figura \ref{fig:fac_cruzada}.

\begin{figure}[!htbp]
    \centering
    \begin{subfigure}[b]{0.48\textwidth}
        \centering
        \includegraphics[width=\linewidth]{Programas/resultados/passo3_facs/fac_startrek_h263_64k_zoom.png}
        \caption{FAC de $X[n]$ (\texttt{h263\_64k}, $0 \le k \le 800$).}
    \end{subfigure}
    \hfill
    \begin{subfigure}[b]{0.48\textwidth}
        \centering
        \includegraphics[width=\linewidth]{Programas/resultados/passo3_facs/fac_startrek_mpeg4_hq_zoom.png}
        \caption{FAC de $Y[n]$ (\texttt{mpeg4\_hq}, $0 \le k \le 800$).}
    \end{subfigure}
    \caption{FACs normalizadas pelo VMQ ($r[0] = 1{,}0$).}
    \label{fig:facs_individuais}
\end{figure}

\begin{figure}[!htbp]
    \centering
    \begin{subfigure}[b]{0.48\textwidth}
        \centering
        \includegraphics[width=\linewidth]{Programas/resultados/passo3_facs/fac_cruzada_zoom.png}
        \caption{FAC cruzada $\rho_{XY}[k]$ (detalhe $0 \le k \le 250$).}
    \end{subfigure}
    \hfill
    \begin{subfigure}[b]{0.48\textwidth}
        \centering
        \includegraphics[width=\linewidth]{Programas/resultados/passo3_facs/fac_cruzada_completa.png}
        \caption{FAC cruzada $\rho_{XY}[k]$ ($0 \le k \le 26.191$).}
    \end{subfigure}
    \caption{Correlação cruzada entre $X[n]$ e $Y_{\text{bloco}}[n]$.}
    \label{fig:fac_cruzada}
\end{figure}

Nos gráficos:
\begin{enumerate}[leftmargin=*, nosep]
    \item \textbf{H.263:} a curva parte de $1{,}0$ em $k = 0$ e decai até o nível DC $m_X^2/\text{VMQ}_X = 0{,}8262$ (Figura \ref{fig:facs_individuais}a).
    \item \textbf{MPEG-4:} a curva oscila com período de $12$ quadros (Figura \ref{fig:facs_individuais}b), por causa dos quadros I do GoP. A linha escura é o ponto médio de cada ciclo de $12$ atrasos, ou seja, passa pelo meio da oscilação e não pelos picos.
    \item \textbf{Correlação cruzada:} fica em torno de zero para todos os atrasos ($\rho_{XY}[0] = -0{,}0048$), ou seja, os dois sinais são praticamente descorrelatados.
\end{enumerate}

% ========================================================
% ITEM 4
% ========================================================
\section{Descorrelação 2-a-2}

\noindent\textbf{Enunciado (Item 4):} \textit{Descorrelatar 2-a-2.}

\medskip
$W_1$ e $W_2$ são descorrelatadas quando a covariância cruzada é nula: $K_{W_1 W_2} = E[W_1 W_2] - m_{W_1}m_{W_2} = 0 \iff \rho_{W_1 W_2} = 0$.

Pelo Capítulo 6 do J.J.W. \cite{JJW}, para o vetor centrado $\underline{\widetilde{Z}}[n] = [X[n] - m_X, \; Y_{\text{bloco}}[n] - m_Y]^T$, a matriz de covariância é:
\begin{equation}
\mathbf{K}_{\underline{Z}} = \begin{bmatrix} \sigma_X^2 & K_{XY}[0] \\ K_{XY}[0] & \sigma_Y^2 \end{bmatrix} = \begin{bmatrix} 254.130{,}08 & -2.963{,}91 \\ -2.963{,}91 & 1.481.147{,}70 \end{bmatrix}
\end{equation}
Define-se $\underline{W}[n] = \mathbf{Q}^T \underline{\widetilde{Z}}[n]$, onde as colunas de $\mathbf{Q}$ são os autovetores de $\mathbf{K}_{\underline{Z}}$. Os autovalores são as variâncias de $W_1$ e $W_2$:
\begin{align}
\lambda_1 &= \sigma_{W_1}^2 = 254.122{,}92\text{ bits}^2, & \lambda_2 &= \sigma_{W_2}^2 = 1.481.154{,}86\text{ bits}^2
\end{align}
e a matriz $\mathbf{Q}$ ($\mathbf{Q}^T \mathbf{Q} = \mathbf{I}$) é:
\begin{equation}
\mathbf{Q} = \begin{bmatrix} -0{,}999997 & -0{,}002416 \\ -0{,}002416 & \phantom{-}0{,}999997 \end{bmatrix}
\end{equation}
A covariância de $\underline{W}[n]$ fica diagonal:
\begin{equation}
\mathbf{K}_{\underline{W}} = \mathbf{Q}^T \mathbf{K}_{\underline{Z}} \mathbf{Q} = \begin{bmatrix} \lambda_1 & 0 \\ 0 & \lambda_2 \end{bmatrix} \implies \rho_{W_1 W_2}[0] = \frac{K_{W_1 W_2}[0]}{\sqrt{\lambda_1 \lambda_2}} = 9{,}27 \times 10^{-19} \approx 0
\end{equation}
Logo $W_1$ e $W_2$ estão descorrelatadas no atraso zero. As Figuras \ref{fig:descorrelacao_cruzada} e \ref{fig:facs_descorrelacionadas} mostram a correlação cruzada e as FACs de $W_1[n]$ e $W_2[n]$.

\begin{figure}[!htbp]
    \centering
    \begin{subfigure}[b]{0.48\textwidth}
        \centering
        \includegraphics[width=\linewidth]{Programas/resultados/passo4_descorrelacao/descorrelacao_cruzada_zoom.png}
        \caption{FAC cruzada $\rho_{W_1 W_2}[k]$ ($0 \le k \le 250$).}
    \end{subfigure}
    \hfill
    \begin{subfigure}[b]{0.48\textwidth}
        \centering
        \includegraphics[width=\linewidth]{Programas/resultados/passo4_descorrelacao/descorrelacao_cruzada_completa.png}
        \caption{Comparação antes e depois da descorrelação.}
    \end{subfigure}
    \caption{Correlação cruzada entre $W_1[n]$ e $W_2[n]$.}
    \label{fig:descorrelacao_cruzada}
\end{figure}

\begin{figure}[!htbp]
    \centering
    \begin{subfigure}[b]{0.48\textwidth}
        \centering
        \includegraphics[width=\linewidth]{Programas/resultados/passo4_descorrelacao/fac_descorrelacionado_w1.png}
        \caption{FAC de $W_1[n]$.}
    \end{subfigure}
    \hfill
    \begin{subfigure}[b]{0.48\textwidth}
        \centering
        \includegraphics[width=\linewidth]{Programas/resultados/passo4_descorrelacao/fac_descorrelacionado_w2.png}
        \caption{FAC de $W_2[n]$.}
    \end{subfigure}
    \caption{FACs de $W_1[n]$ e $W_2[n]$.}
    \label{fig:facs_descorrelacionadas}
\end{figure}

\textbf{Descorrelação não implica independência}, exceto para vetores Gaussianos. Como os sinais não são Gaussianos (Item 1), $W_1[n]$ e $W_2[n]$ são descorrelatados, mas não necessariamente independentes.

% ========================================================
% ITEM 5
% ========================================================
\section{Análise de Estacionariedade (ESA e EE)}

\noindent\textbf{Enunciado (Item 5):} \textit{Os sinais podem ser considerados processos ESA? EE?}

\medskip
Pelo Capítulo 7 do J.J.W. \cite{JJW}:
\begin{itemize}[leftmargin=*, nosep]
    \item \textbf{ESA (Definição 7.10):} média constante ($m(t) = m$) e autocorrelação que depende só da diferença entre os instantes ($R(t_1, t_2) = R(t_2 - t_1)$);
    \item \textbf{EE (Definição 7.9):} todas as f.d.p.s conjuntas são invariantes a deslocamentos no tempo.
\end{itemize}

Para testar essas condições, observa-se o comportamento dos sinais ao longo do tempo (Figura \ref{fig:estacionariedade}).

\begin{figure}[!htbp]
    \centering
    \begin{subfigure}[b]{0.48\textwidth}
        \centering
        \includegraphics[width=\linewidth]{Programas/resultados/passo1_sinais_e_histogramas/sinal_startrek_h263_64k.png}
        \caption{H.263: o nível médio muda ao longo do tempo.}
    \end{subfigure}
    \hfill
    \begin{subfigure}[b]{0.48\textwidth}
        \centering
        \includegraphics[width=\linewidth]{Programas/resultados/passo1_sinais_e_histogramas/sinal_startrek_mpeg4_hq.png}
        \caption{MPEG-4: a dispersão muda ao longo do tempo.}
    \end{subfigure}
    \caption{Sinais ao longo do filme.}
    \label{fig:estacionariedade}
\end{figure}

Conclusão:
\begin{enumerate}[leftmargin=*, nosep]
    \item \textbf{Não são ESA:} a média e a variância mudam ao longo do filme conforme a cena (cenas paradas usam menos bits que cenas de ação). No MPEG-4, a autocorrelação também depende de o quadro ser I, P ou B, e não só do atraso $k$.
    \item \textbf{Não são EE:} todo processo EE com segundo momento finito é ESA. Como os sinais não são ESA, também não são EE.
\end{enumerate}

% ========================================================
% ITEM 6
% ========================================================
\section{Densidade Espectral de Potência (DEP)}

\noindent\textbf{Enunciado (Item 6):} \textit{Plotar as DEPs.}

\medskip
A DEP $S(f)$ é a Transformada de Fourier da autocorrelação (Wiener-Khinchin \cite{JJW}):
\begin{equation}
S(f) = \sum_{k=-\infty}^{\infty} R[k]\,e^{-j 2\pi f k}
\end{equation}
A Figura \ref{fig:deps} mostra as DEPs, calculadas pela FFT da autocorrelação. No MPEG-4 aparecem raias em $f = m/12$, por causa do GoP de $12$ quadros.

\begin{figure}[!htbp]
    \centering
    \begin{subfigure}[b]{0.48\textwidth}
        \centering
        \includegraphics[width=\linewidth]{Programas/resultados/passo6_7_dep_e_espectro/dep_linear_h263.png}
        \caption{DEP de $X[n]$ (\texttt{startrek\_h263\_64k}).}
    \end{subfigure}
    \hfill
    \begin{subfigure}[b]{0.48\textwidth}
        \centering
        \includegraphics[width=\linewidth]{Programas/resultados/passo6_7_dep_e_espectro/dep_linear_mpeg4.png}
        \caption{DEP de $Y[n]$ (\texttt{startrek\_mpeg4\_hq}).}
    \end{subfigure}
    \caption{DEPs em escala linear.}
    \label{fig:deps}
\end{figure}

% ========================================================
% ITEM 7
% ========================================================
\section{Cor do Espectro de Cada Sinal}

\noindent\textbf{Enunciado (Item 7):} \textit{Qual a cor do espectro de cada um dos sinais de vídeo.}

\medskip
A cor do espectro vem da inclinação da DEP em dB por década, nas baixas frequências ($S(f) \propto 1/f^\alpha$):
\begin{itemize}[leftmargin=*, nosep]
    \item \textbf{Ruído Branco ($\alpha = 0$):} DEP plana ($0\text{ dB/década}$);
    \item \textbf{Ruído Rosa ($\alpha = 1$):} queda de $-10\text{ dB/década}$ ($S(f) \propto 1/f$);
    \item \textbf{Ruído Marrom ($\alpha = 2$):} queda de $-20\text{ dB/década}$ ($S(f) \propto 1/f^2$).
\end{itemize}

Medindo a queda da DEP entre $X = 1{,}0$ e $X = 10{,}0$ (uma década), Figura \ref{fig:cor_espectro}:
\begin{itemize}[leftmargin=*, nosep]
    \item \textbf{H.263 (\texttt{startrek\_h263\_64k}):} queda de $10{,}05\text{ dB/década}$, \textbf{Ruído Rosa};
    \item \textbf{MPEG-4 (\texttt{startrek\_mpeg4\_hq}):} queda de $9{,}66\text{ dB/década}$, também \textbf{Ruído Rosa}, com as raias do GoP por cima.
\end{itemize}

\begin{figure}[!htbp]
    \centering
    \begin{subfigure}[b]{0.48\textwidth}
        \centering
        \includegraphics[width=\linewidth]{Programas/resultados/passo6_7_dep_e_espectro/cor_espectro_startrek_h263_64k.png}
        \caption{H.263 (queda de $10{,}05\text{ dB/década}$).}
    \end{subfigure}
    \hfill
    \begin{subfigure}[b]{0.48\textwidth}
        \centering
        \includegraphics[width=\linewidth]{Programas/resultados/passo6_7_dep_e_espectro/cor_espectro_startrek_mpeg4_hq.png}
        \caption{MPEG-4 (queda de $9{,}66\text{ dB/década}$).}
    \end{subfigure}
    \caption{Queda da DEP em dB entre $X = 1{,}0$ e $X = 10{,}0$.}
    \label{fig:cor_espectro}
\end{figure}

% ========================================================
% CONCLUSÃO
% ========================================================
\section*{Conclusão}

Neste trabalho, aplicamos os conceitos de variáveis aleatórias conjuntas e processos estocásticos a sinais reais de vídeo:
\begin{enumerate}[leftmargin=*, nosep]
    \item Os sinais não são Gaussianos: só assumem valores não negativos, o H.263 tem dois picos e o MPEG-4 é assimétrico.
    \item Os momentos confirmam $\text{VMQ} = \sigma^2 + m^2$. O bloco de $26.192$ amostras do MPEG-4 manteve a média e o desvio-padrão da série completa com erro menor que $1\%$.
    \item A FAC do H.263 decai até o nível DC. A do MPEG-4 oscila com período de $12$ quadros (GoP) e a curva central passa pelo meio das oscilações.
    \item A diagonalização de $\mathbf{K}_{\underline{Z}}$ anulou a correlação em $k=0$ ($\rho_{W_1 W_2}[0] \approx 0$). Os sinais não são Gaussianos, então descorrelatados não significa independentes.
    \item Os sinais não são ESA nem EE. As DEPs caem cerca de $10\text{ dB/década}$ ($10{,}05$ no H.263 e $9{,}66$ no MPEG-4), ou seja, o espectro é Ruído Rosa.
\end{enumerate}

% ========================================================
% REFERÊNCIAS BIBLIOGRÁFICAS
% ========================================================
\begin{thebibliography}{9}
\bibitem{JJW} ALBUQUERQUE, J. P. A.; FORTES, J. M. P.; FINAMORE, W. A. \textit{Probabilidade, Variáveis Aleatórias e Processos Estocásticos}. Rio de Janeiro: Editora PUC-Rio / LTC, Capítulos 6 e 7.
\bibitem{Rosangela} COELHO, R. F. \textit{Notas de Aula e Apostilas de Modelos Probabilísticos}. Seção de Engenharia Elétrica e de Sistemas (SE/3), Instituto Militar de Engenharia (IME), 2026.
\end{thebibliography}

\end{document}
